\newpage
{\Large Abstract}\\
\hrule

% Place your abstract here.
% \begin{itemize}
% \item Open the {\bf Project Report Template.tex} file and carefully follow the comments (starting with \%).
% \item Process the file with {\bf pdflatex}, using other processors may need you to change some features such as graphics types.
% \item Note the files included in the  {\bf Project Report Template.tex} (with the .tex extension excluded). You can open these files separately and modify their contents 
% or create new ones.
% \item Contact the latex namual for more features in your document such as equations, subfigures, footnotes, subscripts \& superscripts, special characters etc.
% \item I recommend using the {\bf kile} latex IDE, as it is simple to use.
% \end{itemize}

Teleoperation systems are used to allow humans to control robotic systems over large distances. An ideal user experience involves non-intrusive sensing and intuitive spatial mapping. Few implementations of teleoperation discuss the effects of applying cybersecurity principles to the resource constrained hardware inherent to some robotic systems. This report discusses the implementation of an intuitive and non-invasive teleoperation system, using computer vision, while keeping network security and constraints in consideration. The system uses a pre-trained computer vision pipeline, along with an inverse kinematics algorithm to determine the actuation values for a Tinkerkit Braccio. Packets sent from the operator-side to the Braccio-side are managed by the DTLS protocol, using pre-shared keys. Due to monocular depth ambiguity, the motion of the Braccio was constrained to 2 dimensions. Through experiments, it was found that the smoothing algorithm used to reduce aggressive movements contributed most ($88.3\%$) to the system's total latency. The latency due to encryption and decryption were otherwise insignificant, and this is more apparent over slower networks, such as due to large geographical separation. The implementation code, report source code, raw data and videos can all be found on the project's \href{https://github.com/Neo-Vorsatz-UCT/eee4022s}{GitHub}.
